\clearpage
\section{\label{sec:primer}A Very Brief Primer on Graphing and AI}

\justifying
While an in depth explanation of graph theory is beyond the scope of this paper, some background
will be provided here to make the paper accessible to a wider audience. A collection of related references
will be included at the end of this paper for those who may be interested.

\subsection{\label{sec:Graph}About Graphs}

\justifying
Use of the term graph in this paper simply means a representation of information in a way that describes a group of
things, which we will refer to as ``nodes''. A relationship between nodes is called an ``edge''. Graphs have many
properties, most relevant of which to us include whether a graph has directionality, As with many things in computer
science, we are affected by problems of scale. It often makes sense to constrain our scope to a given sub-graph or ``neighborhood''.
\vspace{2mm}

\justifying
Most likely when this paper refers to a ``graph'' the reader is safe to assume a reference to an image like figure \ref{digraph}. The nodes are represented as colored circles, the edges are represented by arrows. The use of arrows
instead of lines means the graph is ``directional'' (rather than ``undirected'') and will be referred to as a ``digraph''
hereafter.
\vspace{2mm}

\subsection{\label{sec:DL}About Deep Learning with Graph Neural Networks}

\justifying
The first item of note is Consider figure \ref{aimldl}.
\vspace{2mm}

\begin{figure}[H]
	\includegraphics[width=12cm]{rodrigo1-768x332.png}
	\caption{From \href{https://www.ibm.com/blogs/systems/ai-machine-learning-and-deep-learning-whats-the-difference/}{AI, machine learning and deep learning: What’s the difference?} by Rodrigo Ceron.}
	\label{aimldl}
\end{figure}